\section{Experiments on recorded billiards scenes}
\label{sec:experiments}

\subsection{Data, partitions, and measurement}

The source repository contains twenty distinct Real Pool 3D recordings in ten paired scene groups~\citep{projectrepo}. A second directory contains byte-identical video copies; these are excluded from the recording count. The twenty videos contain 2,499 frames at $1280\times720$ resolution and approximately $59.94$ frames per second, totaling $41.69$ seconds. The original observed prefixes contain 3--51 frames. We preserve these prefixes and measure prediction at horizons of 15, 30, and 60 frames where the recording is long enough. One training prefix is already near its first contact; there is no event-oracle segmentation. Prefix lengths and clip counts are properties of this small collection, not independent samples obtained by extracting additional windows.

Groups 1--6 supply training data, groups 7--8 supply validation data, and groups 9--10 supply the four evaluation clips. Both variants of a scene remain in the same partition. An initial diagnostic pass inspected all ten groups before the grouped experiment. Consequently, the evaluation is exploratory: the evaluation clips are excluded from neural training, but are not previously unexamined recordings. We report the two evaluation groups separately and make no statistical-significance or broad generalization claim.

Target positions are measured with a separate image-processing procedure from the Hough detector used for prediction. Non-green foreground components are fitted with ellipses, with local image-template matching when two objects become connected. Initial, prefix-end, and future keyframes, together with candidate contact neighborhoods, are visually checked. Low-contrast cue-ball initialization and a distorted radius receive image-based corrections. Foreground support and image boundaries determine whether a tracked center is measurable. These are image-derived reference measurements, not exact states exported from the game engine or externally validated annotations.

The measured cohort is a subset of trackable circular foreground objects established in the first frame; clipped objects, unseparated dark regions, and missed initial detections are outside that cohort. All measured identities remain in the evaluation record, including unavailable target frames. Prediction tracks, by contrast, are constructed only from the observed prefix, using Hough detections at its endpoint and backward association within the prefix. Unmatched prediction objects incur the image-diagonal penalty in end-to-end displacement scoring. Missing target measurements are excluded from displacement averages and their coverage is reported. The moving-object stratum contains reference objects whose future displacement exceeds five pixels; this is used for reporting only and does not enter inference.

\subsection{Methods and training}

Classical comparisons include frame persistence, a linear fit, quadratic extrapolation, and rigid-body rollouts initialized by the same quadratic state estimate. The rollouts use constant velocity, the fitted constant acceleration, or a speed-opposing deceleration with magnitude equal to the fitted acceleration norm. Deceleration clamps at zero speed to avoid reversal. For contact ablations, shape contacts are disabled while retaining the same initial state, integrator, and acceleration rule. The per-shape material settings are those in Table~\ref{tab:settings}; each video interval contains four physics substeps. Boundaries and pockets are not modeled.

A compact ConvLSTM baseline has one convolutional recurrent layer with 16 hidden channels, a $3\times3$ gate kernel, and a $3\times3$ RGB readout, totaling 11,443 parameters. It predicts at $160\times90$ resolution. Training batches contain four windows drawn only from the twelve training clips, with prefix lengths sampled from $\{3,5,8,11,15,16\}$ and a ten-frame autoregressive rollout. The loss is the mean squared RGB error across that rollout. Adam uses learning rate $10^{-3}$ with gradient-norm clipping at 1.0. Three random seeds are run on an Apple GPU using PyTorch's MPS backend.

The initial budget was 400 updates. Because validation loss was still decreasing, the same architecture and optimizer were rerun with a budget of 2,000 updates per seed. Checkpoints are chosen by the mean 30-frame validation MSE, measured every 100 updates. The earlier budgets and evaluation outcomes are retained with the experiment record. This budget extension is part of the exploratory procedure; it does not supply a new untouched evaluation set. The reported baseline is a small model trained on this dataset, not a reproduction of a published large-scale video-prediction benchmark.

For image comparisons, classical trajectories share a fixed appearance renderer. It takes object patches from the last observed frame, inpaints their original locations, and translates the patches to predicted centers. No future pixel is used to render a prediction. Persistence repeats the complete last observed frame. All rendered and learned outputs are scored against identical $160\times90$ RGB targets. We report complete-frame MSE and mean per-clip PSNR; a largely static background can favor persistence, so image scores are interpreted alongside geometry.

\subsection{Results and contact ablations}

\begin{table}[htbp]
\caption{Thirty-frame geometric prediction on four evaluation clips in groups 9 and 10. Image-derived reference centers define the measured cohort. Means give equal weight to the two paired groups; moving ADE uses reference displacement only for reporting.}
\label{tab:geometry}
\centering
\begin{tabular}{lrrr}
\toprule
Method & ADE (px) & FDE (px) & Moving ADE (px) \\
\midrule
Persistence & 55.46 & 76.91 & 145.63 \\
Linear & 51.70 & 89.51 & 116.59 \\
Quadratic & 41.19 & 65.62 & 90.52 \\
Same state, no contacts & 42.92 & 72.34 & 92.90 \\
Physics (constant velocity) & 12.42 & 21.24 & 28.91 \\
Physics (constant acceleration) & 16.58 & 35.95 & 39.72 \\
Deceleration, no contacts & 39.99 & 62.04 & 88.38 \\
Physics (deceleration) & 14.02 & 26.76 & 34.56 \\
\bottomrule
\end{tabular}
\end{table}

At 30 frames, constant-velocity contact simulation has mean ADE 12.42 pixels compared with 42.92 for its matched contact-disabled rollout and 41.19 for quadratic extrapolation. The corresponding moving-object errors are 28.91, 92.90, and 90.52 pixels. Deceleration gives 14.02 pixels overall and 34.56 on moving objects. Constant-velocity simulation also has lower validation ADE than deceleration; all variants remain reported.
Reference-measurement coverage averages 99.79\% across evaluation clips at the primary horizon, and every reference identity is matched by a prefix track. Pair-time penetration exceeding 10\% of summed radii is 0.00\% for constant-velocity contact simulation versus 2.17\% with contacts disabled. These measurements apply to the annotated object subset.

\begin{table}[htbp]
\caption{Complete RGB reconstruction at 30 frames and $160\times90$ resolution. PSNR is averaged per clip. ConvLSTM variation is the standard deviation across three seeds, not a scene-level confidence interval.}
\label{tab:rgb}
\centering
\begin{tabular}{lrr}
\toprule
Method & RGB MSE & PSNR (dB) \\
\midrule
Persistence & 0.002162 & 26.80 \\
Linear & 0.001953 & 27.37 \\
Quadratic & 0.001804 & 27.62 \\
Same state, no contacts & 0.001699 & 28.21 \\
Physics (constant velocity) & 0.001150 & 29.52 \\
Physics (deceleration) & 0.001428 & 28.53 \\
ConvLSTM (3 seeds) & $0.002727 \pm 0.000226$ & $25.85 \pm 0.30$ \\
\bottomrule
\end{tabular}
\end{table}

The compact ConvLSTM has mean RGB MSE 0.002727 compared with 0.002162 for persistence and 0.001150 for constant-velocity simulation with the shared renderer. Selected checkpoints occur at updates 1,900, 2,000, and 2,000. Its outputs remain blurred in several clips and do not improve on persistence under this image metric. This outcome applies to the stated model and budget, not learned video predictors as a class.

\begin{table}[htbp]
\caption{Per-group 30-frame ADE, averaged over paired clips. Two groups do not support a reliable significance claim.}
\label{tab:pergroup}
\centering
\begin{tabular}{rrrr}
\toprule
Group & Physics (CV) & Same state, no contacts & Quadratic \\
\midrule
9 & 8.17 & 18.43 & 18.09 \\
10 & 16.66 & 67.40 & 64.29 \\
\bottomrule
\end{tabular}
\end{table}



The geometric comparison isolates the effect of modeled interactions within the evaluated cohort. Contact-enabled constant-velocity simulation is compared directly with its contact-disabled counterpart, using the same estimated states. Deceleration and constant acceleration are also reported, rather than selecting an acceleration rule from evaluation-scene outcomes. Candidate near-contact timing and object-pair counts provide descriptive diagnostics, but proximity alone can count a pass-through as contact. Displacement and persistent penetration are therefore needed to assess the consequences of an interaction.

Synchronized replays of all four evaluation clips are available at \url{https://aabbccdkg.github.io/Video_prediction_through_physical_laws/}. The page compares contact ablations and all three ConvLSTM seeds, with frame stepping, input/future boundaries, slower contact playback, and links to the paper and experiment files. Display enlargement preserves the $160\times90$ prediction resolution.

The image comparison evaluates the complete prediction, including appearance. The compact ConvLSTM's outputs and training variability are included even where they fail to improve on persistence. Its blurred objects make automatic extraction of centers unreliable; detector-dependent neural displacement penalties remain in the experiment record rather than being presented as a directly comparable primary trajectory score. Figure~\ref{fig:qualitative} shows every evaluated clip at the primary horizon with a fixed seed for the neural illustration.

